% =============================================================================
% Seção 6 — Conclusão
% =============================================================================

\section{Conclusão}
\label{sec:conclusao}

Apresentamos o NeuroAssist AI, uma plataforma de comunicação assistiva
offline-first e preservadora de privacidade que transforma gestos faciais residuais
de pacientes com esclerose lateral amiotrófica em estágio avançado em alertas de
emergência estruturados, entregues a cuidadores por uma rede local usando hardware
comum.

As três contribuições originais do sistema abordam lacunas distintas, mas
inter-relacionadas, na tecnologia assistiva atual para essa população. Primeiro, ao
executar a detecção de landmarks faciais inteiramente dentro de uma sandbox
WebAssembly e servindo o modelo binário como ativo local, o NeuroAssist AI alcança
localidade estrita de dados de vídeo sem sacrificar a precisão ou a taxa de quadros
de soluções server-side. Segundo, a máquina de estados de gesto deliberado de seis
estados — que exige fechamento de duração mínima, reabertura confirmada entre piscadas
e uma janela temporal — reduz substancialmente os falsos positivos em relação a
comparações simples de limiar EAR por quadro, como quantificado na
Seção~\ref{sec:avaliacao} (taxa de falsos positivos: [B]\,\% vs.\ [A]\,\% para o
limiar simples; $N$=[N] participantes, [M]\,min cada). Terceiro, a arquitetura de
microsserviços hexagonais com aplicação ArchUnit em tempo de compilação e o mecanismo
plug-in \texttt{engineFeatureRef} fornece uma base extensível que desacopla a detecção
de gesto do gerenciamento clínico de comandos, facilitando a expansão clínica
incremental sem regressão arquitetural.

A latência ponta-a-ponta da conclusão do gesto até a exibição do alarme ao cuidador
foi em média [X]\,ms ($\sigma$=[Y]\,ms, $N$=[N] ensaios) em uma rede local com o
transporte de polling de cinco segundos. O perfil de latência é clinicamente aceitável
para sinalização de emergência: o componente dominante é o intervalo de polling
($\leq$5\,s), e não o tempo de detecção ou processamento no servidor, e a substituição
por eventos enviados pelo servidor reduziria a latência total para aproximadamente
[X]\,ms.

O custo de hardware (webcam comum, notebook ou computador de placa única padrão, sem
GPU) posiciona o NeuroAssist AI como uma opção viável para implantação domiciliar em
países de baixa e média renda, onde reside a maioria da população global com ELA e
onde dispositivos especializados de CAA são economicamente inacessíveis. Juntamente
com a arquitetura de privacidade desde a concepção, que satisfaz os requisitos da
LGPD para dados biométricos de saúde sem infraestrutura institucional de processamento
de dados, essas propriedades abordam as barreiras que impedem os sistemas comerciais
atuais de CAA de alcançar pacientes em cuidado domiciliar.

A principal limitação deste trabalho é a escala da avaliação: o estudo de caso envolve
um paciente com ELA e cinco cuidadores e voluntários, e seus achados não devem ser
generalizados para a população mais ampla com ELA sem um ensaio multicêntrico sob
aprovação de comitê de ética em pesquisa (CEP). O protocolo de calibração
específico por paciente e a infraestrutura de medição produzidos neste trabalho são
projetados para apoiar esse ensaio com adaptação mínima, e os resultados preliminares
do estudo de caso fornecem estimativas concretas de latência, taxa de falsos positivos
e recall de detecção para calcular o tamanho de amostra de estudos futuros.
